\documentclass[11pt]{article}
\usepackage{preamble}
\setlength{\parskip}{4pt}

\begin{document}

\daybanner{AP Statistics \textperiodcentered\ Topic 5.4 (CED 3.1) \textperiodcentered\ B: 2026-11-20 \textperiodcentered\ E: 2026-12-18 \textperiodcentered\ TEACHER KEY}{Centered or Close?}

\sectionbanner{CED TETHER}

{\small
\begin{itemize}[leftmargin=1.5em,itemsep=2pt,topsep=2pt]
  \item Enduring Understanding: UNC-3 Probabilistic reasoning allows us to anticipate patterns in data.
  \item Skills:
  \item \textbf{Skill 4.B} --- Interpret statistical calculations and findings to assign meaning or assess a claim.
  \item \textbf{Skill 3.B} --- Determine parameters for probability distributions.
  \item Learning Objectives and Essential Knowledge:
  \item LO UNC-3.I Explain why an estimator is or is not unbiased. [Skill 4.B]
  \item EK UNC-3.I.1 When estimating a population parameter, an estimator is unbiased if, on average, the value of the estimator is equal to the population parameter.
  \item LO UNC-3.J Calculate estimates for a population parameter. [Skill 3.B]
  \item EK UNC-3.J.1 When estimating a population parameter, an estimator exhibits variability that can be modeled using probability.
  \item EK UNC-3.J.2 A sample statistic is a point estimator of the corresponding population parameter.
\end{itemize}
}
\textbf{Essential question:} How should purpose change the way we weigh an estimator's bias and variability?\par
\textbf{DOK-3 skill:} 4.B \textperiodcentered\ \textbf{FRQ pattern:} {\small\texttt{unbiased-\allowbreak{}vs-\allowbreak{}low-\allowbreak{}variability-\allowbreak{}which-\allowbreak{}estimator-\allowbreak{}for-\allowbreak{}the-\allowbreak{}purpose}}\par
\textbf{DOK rationale:} Part (c) weighs long-run center, spread, and purpose-specific accuracy against two absolute claims, chooses an estimator for the stated loss, and explains the consequence of ignoring its bias.\par

\frameworkphaseheader{Do Now (first take) $\rightarrow$ Explore (video + follow-along) $\rightarrow$ Exit (finish + turn in)}{1 $\rightarrow$ 3}{45}{%
\begin{itemize}[leftmargin=1.5em,itemsep=2pt,topsep=2pt]
  \item Board slide up at the bell. Record Lena, Omar, and depends votes without endorsing a universal rule.
  \item Begin the Topic 5.4 video follow-along at minute 5; leave first takes ungraded.
  \item Return to the board slide at minute 33. Circulate and separate one estimate from a sampling distribution.
  \item Collect at the door. Score part (c) only, E/P/I.
\end{itemize}
}{%
\begin{itemize}[leftmargin=1.5em,itemsep=2pt,topsep=2pt]
  \item Choose a claim and cite one table entry before the lesson.
  \item Complete the follow-along about point estimators, bias, and sampling variability.
  \item Finish (a)--(c), revise the first take in the exit line, and turn in the sheet.
\end{itemize}
}{%
\begin{itemize}[leftmargin=1.5em,itemsep=2pt,topsep=2pt]
  \item Does Rule B landing exactly on 0.40 once tell you where it centers over all samples?
  \item Which table column measures systematic direction, and which measures sample-to-sample spread?
  \item What does the phrase within 0.06 make the team's loss function reward?
  \item How could a 2.9-point upward bias matter in a threshold decision?
\end{itemize}
}{Solo teacher. Circulate ~5 pairs. Require students to name the stated goal before declaring an estimator better and to critique the word always in both claims.}

\begin{callout}{\IconWarn\ WATCH FOR}{calloutred}
\begin{itemize}[leftmargin=1.5em,itemsep=2pt,topsep=2pt]
  \item Using the $X=9$ sample to classify long-run bias.
  \item Treating smaller standard deviation as proof that Rule B is unbiased.
  \item Choosing the unbiased rule without using the within-0.06 criterion.
\end{itemize}
\end{callout}

\sectionbanner{SCORING PART (c) --- E / P / I}

\textbf{Expected elements}
\begin{itemize}[leftmargin=1.5em,itemsep=2pt,topsep=2pt]
  \item \textbf{separates-bias-from-variability} (required) --- Identifies Rule A as centered at 0.40 and Rule B as upward biased while recognizing Rule B's smaller standard deviation
  \item \textbf{chooses-for-stated-purpose} (required) --- Chooses Rule B for the within-0.06 goal and cites probabilities 0.659 versus 0.459
  \item \textbf{rejects-always-language} (required) --- Explains that unbiasedness does not guarantee a close single estimate and lower variability does not erase systematic bias
  \item \textbf{states-reader-consequence} (required) --- Explains that Rule B can systematically make 40 percent support look like about 43 percent or influence a threshold decision
\end{itemize}

{\small\renewcommand{\arraystretch}{1.25}
\noindent\begin{tabular}{|p{0.5in}|p{5.9in}|}
\hline
\textbf{E} & Correctly distinguishes center from spread, chooses Rule B for the stated accuracy goal with the numerical comparison, rejects both absolute claims, and gives a concrete consequence of Rule B's upward bias. \\ \hline
\textbf{P} & Makes the purpose-supported choice and identifies the bias-variability tradeoff but omits one numerical comparison, one claim critique, or the reader consequence. \\ \hline
\textbf{I} & Calls Rule B unbiased because one estimate equals 0.40, assumes unbiased means every estimate is accurate, or chooses without relating the table to the stated goal. \\ \hline
\end{tabular}}

\textbf{Common mistakes}
\begin{itemize}[leftmargin=1.5em,itemsep=2pt,topsep=2pt]
  \item Judging bias from the single sample with $X=9$ instead of the long-run mean
  \item Comparing 0.070 and 0.098 as centers rather than standard deviations
  \item Calling Rule A always better solely because it is unbiased
  \item Calling Rule B unbiased because its bias is smaller than the 0.06 tolerance
\end{itemize}

\clearpage
\focusbanner{TODAY'S PROBLEM --- annotated}

Suppose a research team tests two rules for estimating a support proportion in a validation population where the true proportion is $p=0.40$. For a random sample of 25 people, let $X$ be the number who support the proposal.\par\smallskip Rule A reports $\widehat p_A=X/25$. Rule B reports $\widehat p_B=(X+5)/35$. Exact summaries of the two sampling distributions are shown.\par\smallskip \textbf{Lena:} ``A is unbiased, so it is always the more trustworthy estimator.''\par \textbf{Omar:} ``B varies less and is more likely to land within 0.06 of the truth, so its bias does not matter.''\par
\begin{center}
\textbf{Sampling distributions when $p=0.40$}\par\smallskip
\begin{tabular}{|l|c|c|c|}
\hline
\textbf{Rule} & \textbf{Mean} & \textbf{SD} & \textbf{Within 0.06} \\ \hline
\textbf{A} & 0.400 & 0.098 & 0.459 \\ \hline
\textbf{B} & 0.429 & 0.070 & 0.659 \\ \hline
\end{tabular}
\end{center}
\begin{firsttakebox}{FIRST TAKE (5 min)}
Whose recommendation do you trust more, Lena's or Omar's? Cite one entry from the table.\par
\answer{Not graded. Each claim names a real strength, then overgeneralizes it.}
\end{firsttakebox}

\textit{Rules callout ``BIAS AND VARIABILITY ANSWER DIFFERENT QUESTIONS (keep on the page)'' is printed on the student sheet.}\par\medskip

\dokbadge{1}~\textbf{(a)} If one sample has $X=9$ supporters, calculate the point estimate from Rule A and from Rule B. State the population parameter both rules estimate.\par
\answer{For $X=9$, Rule A gives $9/25=0.36$ and Rule B gives $(9+5)/35=14/35=0.40$. Both statistics are point estimators of the population support proportion $p$.}

\dokbadge{2}~\textbf{(b)} Use the table to classify each rule as biased or unbiased. For any biased rule, give the direction and amount of bias; then compare the rules' sampling variability.\par
\answer{Rule A is unbiased because its long-run mean, 0.400, equals $p=0.40$. Rule B is biased upward because its long-run mean is $15/35=0.428571\ldots$, so its bias is $0.428571-0.40=0.028571$, about 2.9 percentage points. Rule B is less variable because its standard deviation is about 0.070 compared with about 0.098 for Rule A.}

\dokbadge{3}~\textbf{(c)} Adjudicate both claims. For the goal of maximizing the chance that one estimate is within 0.06 of $p$, choose a rule using its center, spread, and within-0.06 probability. Explain why neither criterion means always better and what ignoring Rule B's bias could mislead a reader into believing.\par
\answer{For the within-0.06 goal, use Rule B: its success probability is 0.659 versus 0.459 for Rule A, and its smaller SD, 0.070 versus 0.098, offsets its 0.029 upward bias for this tolerance. Lena is wrong that unbiased means closer in one sample. Omar's choice fits this goal, but bias still matters: Rule B centers near 0.429 rather than 0.400. Ignoring that bias could make 40 percent support look like 43 percent and influence a threshold decision; another cost or goal could favor Rule A.}

\begin{summaryexitbox}{EXIT}
One thing the video changed about my first take: \hrulefill
\end{summaryexitbox}

\end{document}
